% =====================================================================
% SYNTHETIC SAMPLE DRAFT (OVER-LENGTH) — for demonstrating the
% academic-paper-revision skill's page-fit / trimming step (Step 4).
% Fictional paper: authors, dataset, and numbers are invented, no IP.
% Target venue for the demo: an ACL-family SHORT paper (4-page body limit,
% excluding References / Limitations / Ethics). This draft intentionally runs
% ~5-6 pages so the skill must trim. The trimming levers are deliberate:
%   - two overlapping Introduction paragraphs (redundant motivation)
%   - a verbose, repetitive Related Work
%   - a long inline WORKED EXAMPLE that belongs in an appendix
%   - Discussion paragraphs that restate the results
%   - hedgy filler sentences throughout
% Unlike the student-draft sample, this one is otherwise fairly complete:
% prose contributions (no RQ scaffolding), one equation, cited .bib.
% =====================================================================
\documentclass[11pt]{article}

\usepackage[utf8]{inputenc}
\usepackage{graphicx}
\usepackage{booktabs}
\usepackage{amsmath}
\usepackage{url}

\title{A Rule-Based Detector for Passive Voice in Academic Abstracts}
\author{Alex Sample \and Robin Placeholder \\ School of Made-Up Studies, Institute of Nowhere \\ \texttt{\{asample, rplaceholder\}@example.edu}}
\date{}

\begin{document}
\maketitle

\begin{abstract}
Passive voice is common in academic writing and is often discouraged by style
guides. We present a lightweight, rule-based detector that flags passive-voice
clauses in English abstracts using part-of-speech patterns over auxiliary and
participle sequences. On a hand-labelled set of 600 sentences the detector
reaches an F1 of 0.91, competitive with a supervised classifier that requires
training data, while remaining fully interpretable and dependency-free. We
release the rule set and the annotation guidelines. The method is intended as a
transparent writing-support component rather than a state-of-the-art parser.
\end{abstract}

\section{Introduction}

Academic style guides frequently advise authors to prefer active over passive
voice, on the grounds that the active voice is more direct and easier to read.
Tools that can automatically flag passive constructions are therefore useful in
writing-support settings, where an author wants quick, interpretable feedback on
a draft. Most existing writing assistants either rely on opaque machine-learning
models or on shallow heuristics that miss many constructions, and there is a gap
for a method that is both transparent and reasonably accurate.

Automatically detecting passive voice is helpful because writers often want fast
feedback while drafting, and an interpretable signal is easier to trust and act
on than an opaque score. Many current writing tools use black-box models or very
shallow rules, so a method that is simultaneously transparent and accurate would
fill a real gap. In this paper we pursue exactly that goal. (This paragraph and
the one above make the same point twice and should be merged.)

We present a rule-based detector that operates over part-of-speech tags and a
small set of clause-level patterns. Our contributions are: (i) a compact,
documented rule set for English passive-voice detection that needs no training
data; (ii) a hand-labelled evaluation set of 600 sentences with annotation
guidelines; and (iii) an evaluation showing the rules are competitive with a
supervised baseline while remaining fully interpretable.

\section{Related Work}

Passive-voice detection has been approached in several ways over the years. Early
grammar checkers used simple pattern matching on auxiliary verbs. Later work
built supervised classifiers on top of syntactic features, which improved recall
but required labelled data and a parser. More recent systems use neural sequence
labellers, which perform well but are opaque and data-hungry. There is also a
large body of work on readability and writing quality more generally, which
sometimes includes passive voice as one of many features. Style guides such as
common manuals of academic writing recommend limiting passive voice, and several
commercial writing assistants implement passive detection as a feature, although
their internal methods are usually undocumented. Our approach differs from the
supervised and neural lines of work in that it needs no training data and is
fully interpretable, and it differs from the shallow grammar-checker heuristics
in that it covers a broader set of clause patterns. In short, prior work is
either accurate but opaque and data-hungry, or transparent but shallow; we aim
for the transparent-and-accurate middle. (This section repeats the
same accurate-vs-transparent contrast three times and can be tightened to two
short paragraphs.)

\section{Method}

Given a sentence, we run a part-of-speech tagger and then scan for a passive
pattern: a form of the auxiliary \emph{to be} (or \emph{to get}) followed,
within a short window, by a past participle, optionally preceded by adverbs and
optionally followed by a \emph{by}-phrase. We score each candidate clause and
flag the sentence as passive if any clause scores above a threshold. Formally,
for a clause $x$ with auxiliary indicator $\mathrm{aux}(x)$ and participle
indicator $\mathrm{part}(x)$ within window $w$, the clause score is

\begin{equation}
    \mathrm{score}(x) = \mathrm{aux}(x)\cdot \mathrm{part}(x, w)\cdot
    \big(1 + \lambda\,\mathrm{by}(x)\big),
    \label{eq:score}
\end{equation}

where $\mathrm{by}(x)\in\{0,1\}$ marks a following \emph{by}-agent and $\lambda$
up-weights clauses with an explicit agent. The sentence is flagged when
$\max_x \mathrm{score}(x)$ exceeds a fixed threshold, chosen on a small
development set separate from the test set.

\paragraph{Worked example.}
Consider the sentence ``The results were carefully analysed by three independent
reviewers before the final decision was made.'' The tagger yields the sequence
determiner, noun, auxiliary (\emph{were}), adverb (\emph{carefully}), past
participle (\emph{analysed}), preposition (\emph{by}), and so on. The first
clause matches the pattern: \emph{were} is an auxiliary form of \emph{to be},
\emph{analysed} is a past participle within the window, \emph{carefully} is an
intervening adverb that the pattern tolerates, and \emph{by three independent
reviewers} supplies a \emph{by}-agent, so $\mathrm{by}(x)=1$ and the clause
receives the agent up-weight. The second clause, ``the final decision was
made,'' also matches: \emph{was} is the auxiliary and \emph{made} the participle,
with no \emph{by}-agent, so $\mathrm{by}(x)=0$. Because at least one clause
exceeds the threshold, the whole sentence is flagged as passive, and because two
clauses match, the detector can additionally report a passive-clause count of
two. This step-by-step walk-through is useful for readers new to the pattern but
is long; in a 4-page short paper it should move to an appendix with only a
one-sentence pointer left in the body.

\section{Experiments}

We built an evaluation set of 600 English sentences sampled from abstracts across
six fields, and two annotators labelled each sentence as active or passive with
an inter-annotator agreement of $\kappa = 0.88$; disagreements were adjudicated.
We compare our rules against a supervised logistic-regression classifier trained
on 2{,}000 separately labelled sentences with part-of-speech n-gram features.

\begin{table}[t]
    \centering
    \begin{tabular}{lccc}
        \toprule
        Method & Precision & Recall & F1 \\
        \midrule
        Shallow aux-only heuristic & 0.94 & 0.72 & 0.82 \\
        Supervised LR (needs training) & 0.90 & 0.93 & 0.92 \\
        Ours (rules, no training) & 0.92 & 0.90 & 0.91 \\
        \bottomrule
    \end{tabular}
    \caption{Passive-voice detection on 600 hand-labelled sentences. Intervals
    omitted for space.}
    \label{tab:results}
\end{table}

Table~\ref{tab:results} shows that our rule set reaches F1 0.91, within one point
of the supervised classifier and well above the aux-only heuristic, while
requiring no training data. The gap to the supervised model comes mostly from
recall on reduced passives without an overt auxiliary, which our patterns do not
cover.

\section{Discussion}

Our results indicate that a compact, interpretable rule set can match a
supervised classifier on this task while avoiding any training-data requirement.
This matters for writing-support tools, where transparency and the ability to
explain a flag to the author are valuable. The main errors are reduced passives
and a handful of tagger mistakes on rare participles.

As the numbers show, the rule-based detector performs almost as well as the
supervised model without needing labelled training data, which is convenient for
deployment and makes the flags easy to explain. Reduced passives and occasional
tagging errors account for most of the remaining mistakes. (This paragraph
restates the one above and should be merged.)

\section{Conclusion}

We presented a transparent, training-free rule set for passive-voice detection in
academic abstracts that is competitive with a supervised baseline. Because it is
interpretable and dependency-light, it is well suited to writing-support tools.
Future work includes covering reduced passives and extending the rules to other
languages.

\begin{thebibliography}{9}
\bibitem{manning2014} Manning, C. et al. (2014). The Stanford CoreNLP Natural Language Processing Toolkit.
\bibitem{honnibal2020} Honnibal, M. et al. (2020). spaCy: Industrial-Strength Natural Language Processing.
\bibitem{biber1999} Biber, D. et al. (1999). Longman Grammar of Spoken and Written English.
\end{thebibliography}

\end{document}
